\section*{Discussion}

This paper establishes that choosing which representation pairing to evaluate --- cheaply, from
properties measured online rather than modelled in advance --- dominates the return available from
accelerating any single kernel. A strategy built around avoiding work wins essentially every
asymmetric measurement taken, and the resulting selector outperforms a competitively tuned Roaring
baseline on every real corpus tested, with no single pairing winning enough of the corpus set to make
the matrix itself unnecessary. The shape of the contribution mirrors Roaring's own: array containers,
bitset containers, run containers and SIMD intersection were each individually known before Roaring
existed, and what made the format matter in practice was their composition into one adaptively
dispatched system, evaluated as a whole rather than argued for container by container
\cite{chambi2016roaring,lemire2018roaring}. The claim advanced here is structured the same way --- a
selector over a wider pairing matrix, gated on measured rather than modelled properties, is the
contribution, not any cell inside it.

Roaring already performs adaptive, per-chunk representation choice, so the live question is not
whether adaptivity helps but whether one fixed choice --- a constant chunk width, a constant
cardinality threshold, set once at ingest --- can be right across the whole density axis a corpus is
later evaluated at. This paper's own analysis says it cannot: a fixed threshold is tested against a
chunk's own cardinality rather than the corpus's global density, so a corpus can be dense in aggregate
while every individual chunk stays under the threshold (Table~\ref{tab:census}). Repairing that
threshold inside the incumbent shifts the advantage in opposite directions at the two ends of the
density axis, so neither constant is right and no single fixed threshold serves both regimes. The fix
is offered here as a contribution to the incumbent, not a criticism of it --- Roaring's default is a
memory-sizing crossover, and it is exactly correct for the objective it was chosen for, which is not
this one. What matters for the argument is that the work-avoidance layer keeps its advantage against
the repaired baseline, because that advantage never depended on the kernel underneath being slow.
Roaring needs adaptive selection across density for the same reason this paper's own selector does;
that is the same argument, tested a second time, on someone else's code.

The mechanism behind the asymmetric-cell result is structural rather than incidental. Those cells are
already work-proportional to the smaller operand, so a wider vector instruction has little left to
accelerate; a rank lookup that answers an entire run without reading its bits, or a zone map that
proves two rows disjoint from a fraction of their footprint, wins by removing work rather than by
processing it faster (Fig.~\ref{fig:strategy-decay}a). The complement carries equal weight:
$\cell{\rB}{\rB}$ is where work is genuinely irreducible in the generic case, and it is exactly there
that two decades of vectorized popcount work \cite{mula2018popcount} correctly pay off, including in
this paper's own dense-cell kernel. Avoiding work and accelerating it are complementary axes, not
competing ones, each correct on its own side of the boundary the density sweep locates. That boundary
also responds to the memory hierarchy in opposite directions: vector throughput is a per-cycle
constant a wider core keeps buying more of, realized fully only while data is fed at that rate,
whereas work avoidance removes a count of memory-hierarchy crossings that widening a core does not by
itself reduce (Fig.~\ref{fig:strategy-decay}b). \spec{If core width and the ratio of cache capacity to
DRAM bandwidth continue to diverge as recent processor generations suggest, the boundary measured
here should move further toward avoidance rather than stay where it is} --- a prediction about future
hardware, not a claim about the hardware tested in this study.

The property doing the work throughout this paper is heterogeneity, not sparsity. A corpus in which
every row sits at the same density has nothing for a pairing matrix to exploit, however sparse or
dense that shared density is; what representation selection is paid for is variance in density across
rows that share one universe and one cost model, of which sparsity is only the more familiar half ---
the dense tail rewards exactly the same reasoning, mirrored around the complement. \spec{Read this
way, the same argument should transfer to any workload that mixes elements of very different density
under a shared representation and a shared cost model: sparse linear-algebra format selection among
dense, compressed-sparse-row and coordinate tiles, or a columnar query engine deciding which blocks a
predicate can skip without scanning, are settings where the condition plausibly holds.} Neither has
been tested here, and the claim is offered as where the mechanism should generalize, not as a result
obtained in either domain.

Several things remain open, stated plainly. The selector is cheap relative to the saving it protects,
but it is not close to optimal: nothing measured here closes the residual gap to a bucket oracle that
already knows the winning cell (Fig.~\ref{fig:selection-cost}b). Every measurement reported here is
single-threaded; the real-corpus comparison ran on \nm{one microarchitecture} while the synthetic
cross-architecture comparison spans \nm{four}; and the tile pipeline has so far been evaluated on
within-tile pairs only, with the cross-tile case and the transpose-build cost left to the end-to-end
evaluation that follows this paper. Most directly, the pairing matrix, the selection gate and the tile
pipeline described here are not currently reachable through the project's public C ABI: this paper
reports a systematic measurement of the underlying mechanism, not a released tool, and the ABI and
packaging work needed to close that gap is sequenced immediately after it, with a genome-wide
linkage-disequilibrium application, Tomahawk \cite{tomahawk_github}, as the eventual, and strictly
later, use of the mechanism rather than its motivation. What should outlast any single ratio in this
paper is the map itself --- which representation pairing wins where, and why --- because that map is
transferable independent of the number attached to any one cell.
